  \documentclass{ximera}

  \input{../../../preamble.tex}
  

\definecolor{fvdnavy}{HTML}{071D43}
\definecolor{fvdcyan}{HTML}{20BFC6}
\definecolor{fvdcyanlight}{HTML}{EFFCFB}
\definecolor{fvdmagenta}{HTML}{F10D69}
\definecolor{fvdmagentalight}{HTML}{FFF1F7}
\definecolor{fvdtext}{HTML}{163044}





%=================================================
% Pedagogische omgevingen
% PDF: tcolorbox
% HTML: learning-box
%=================================================

\newenvironment{voorkennis}
{%
    \ifdefined\HCode
        \HCode{%
<section class="learning-box learning-box--prior">
<div class="learning-box__header">
<span class="learning-box__icon" aria-hidden="true"></span>
<span class="learning-box__title">Voorkennis</span>
</div>
<div class="learning-box__content">
        }%
        \begin{itemize}
    \else
       \begin{tcolorbox}[
    enhanced,
    breakable,
    colback=fvdcyanlight,
    colframe=fvdcyan,
    coltext=fvdtext,
    boxrule=0.5pt,
    leftrule=4pt,
    arc=4mm,
    outer arc=4mm,
    left=7mm,
    right=6mm,
    top=5mm,
    bottom=5mm,
    before skip=6mm,
    after skip=6mm,
    title=Voorkennis,
    coltitle=fvdcyan,
    fonttitle=\bfseries\Large,
    attach title to upper={\par\smallskip},
    boxed title style={
        empty,
        left=0mm,
        right=0mm,
        top=0mm,
        bottom=0mm
    },
    overlay={
        \draw[
            fvdcyan,
            line width=0.5pt
        ]
        ([xshift=42mm,yshift=-13mm]frame.north west)
        --
        ([xshift=-7mm,yshift=-13mm]frame.north east);
    }
]
        \begin{itemize}
    \fi
}
{%
    \end{itemize}
    \ifdefined\HCode
        \HCode{%
</div>
</section>
        }%
    \else
        \end{tcolorbox}
    \fi
}


\newenvironment{leerdoelen}
{%
    \ifdefined\HCode
        \HCode{%
<section class="learning-box learning-box--goals">
<div class="learning-box__header">
<span class="learning-box__icon" aria-hidden="true"></span>
<span class="learning-box__title">Leerdoelen</span>
</div>
<div class="learning-box__content">
        }%
        \begin{itemize}
    \else
\begin{tcolorbox}[
    enhanced,
    breakable,
    colback=fvdmagentalight,
    colframe=fvdmagenta,
    coltext=fvdtext,
    boxrule=0.5pt,
    leftrule=4pt,
    arc=4mm,
    outer arc=4mm,
    left=7mm,
    right=6mm,
    top=5mm,
    bottom=5mm,
    before skip=6mm,
    after skip=6mm,
    title=Leerdoelen,
    coltitle=fvdmagenta,
    fonttitle=\bfseries\Large,
    attach title to upper={\par\smallskip},
    boxed title style={
        empty,
        left=0mm,
        right=0mm,
        top=0mm,
        bottom=0mm
    },
    overlay={
        \draw[
            fvdmagenta,
            line width=0.5pt
        ]
        ([xshift=42mm,yshift=-13mm]frame.north west)
        --
        ([xshift=-7mm,yshift=-13mm]frame.north east);
    }
]
        \begin{itemize}
    \fi
}
{%
    \end{itemize}
    \ifdefined\HCode
        \HCode{%
</div>
</section>
        }%
    \else
        \end{tcolorbox}
    \fi
}


  \addPrintStyle{../../..}

  \title{Uitspraken verbinden}
  \author{Ingmar Herreman}

  \begin{abstract}
  In het vorige hoofdstuk maakten we kennis met wiskundige uitspraken:
  beweringen waarvan we kunnen bepalen of ze waar of onwaar zijn.

  Maar wiskundige redeneringen bestaan zelden uit slechts één uitspraak.
  We combineren uitspraken voortdurend met woorden zoals
  \textbf{niet}, \textbf{en} en \textbf{of}.

  In dit hoofdstuk onderzoeken we hoe zulke samengestelde uitspraken
  worden opgebouwd. We leren de logische operatoren negatie,
  conjunctie en disjunctie kennen en onderzoeken met waarheidstabellen
  hoe de waarheidswaarde van een samengestelde uitspraak afhangt van
  de waarheidswaarden van haar onderdelen.

  Daarbij besteden we bijzondere aandacht aan het verschil tussen
  dagelijkse taal en wiskundige logica. Vooral het woord \emph{of}
  heeft in de wiskunde een precieze betekenis die soms verschilt van
  het gewone taalgebruik.

  Zo bouwen we stap voor stap de formele taal op waarmee we later
  complexere wiskundige uitspraken, implicaties en redeneringen
  kunnen analyseren.
  \end{abstract}

  \begin{document}

  % FVD_CHAPTER_DATA_START
% Automatisch gegenereerd uit index.tex.
% Niet handmatig aanpassen.
\fvdchapterdata{1}{Logica — De taal van correct redeneren}{2}
% FVD_CHAPTER_DATA_END













  \maketitle


  % ============================================================
  % VOORKENNIS
  % ============================================================

  \begin{voorkennis}
    \item Je weet wat een wiskundige uitspraak is.
    \item Je kan de waarheidswaarde van een eenvoudige uitspraak bepalen.
    \item Je weet wat een open uitspraak en een predicaat zijn.
    \item Je begrijpt dat het gekozen domein de betekenis van een uitspraak kan beïnvloeden.
    \item Je kan uitspraken benoemen met letters zoals $p$, $q$ en $r$.
  \end{voorkennis}


  % ============================================================
  % LEERDOELEN
  % ============================================================

  \begin{leerdoelen}
    \item Je kan de negatie van een uitspraak interpreteren en noteren.
    \item Je kan bepalen wanneer een negatie waar of onwaar is.
    \item Je kan twee uitspraken verbinden met een conjunctie.
    \item Je kan twee uitspraken verbinden met een disjunctie.
    \item Je kan de symbolen $\neg$, $\land$ en $\lor$ correct gebruiken.
    \item Je kan de waarheidswaarde van samengestelde uitspraken bepalen.
    \item Je kan waarheidstabellen opstellen en interpreteren.
    \item Je begrijpt dat het wiskundige ``of'' inclusief is.
    \item Je kan gewone taal vertalen naar logische symbolen en omgekeerd.
    \item Je kan samengestelde uitspraken analyseren en hun logische structuur herkennen.
  \end{leerdoelen}


  % ============================================================
  \FVDsection{Van eenvoudige naar samengestelde uitspraken}
  % ============================================================

  Beschouw de uitspraken

  \[
  p:\quad 12\text{ is even}
  \]

  en

  \[
  q:\quad 12\text{ is deelbaar door }3.
  \]

  Beide uitspraken zijn waar.

  We kunnen uit deze twee uitspraken nieuwe uitspraken vormen, bijvoorbeeld:

  \begin{center}
  $12$ is even \textbf{en} $12$ is deelbaar door $3$.
  \end{center}

  of

  \begin{center}
  $12$ is even \textbf{of} $12$ is deelbaar door $3$.
  \end{center}

  We kunnen ook één uitspraak ontkennen:

  \begin{center}
  $12$ is \textbf{niet} even.
  \end{center}

  De woorden

  \[
  \text{niet},\qquad \text{en},\qquad \text{of}
  \]

  hebben in de logica een precieze betekenis.

  Ze worden \textbf{logische operatoren} genoemd.


  \begin{definitie}
  Een \textbf{samengestelde uitspraak} is een uitspraak die wordt opgebouwd
  uit één of meer andere uitspraken met behulp van logische operatoren.
  \end{definitie}


  In dit hoofdstuk bestuderen we drie fundamentele operatoren:

  \[
  \neg \qquad \land \qquad \lor
  \]

  met respectievelijk de betekenis

  \[
  \text{niet},\qquad \text{en},\qquad \text{of}.
  \]


  % ============================================================
  \FVDsection{De negatie}
  % ============================================================

  We beginnen met de eenvoudigste logische bewerking.

  Beschouw

  \[
  p:\quad 15\text{ is deelbaar door }3.
  \]

  Deze uitspraak is waar.

  De uitspraak

  \begin{center}
  $15$ is niet deelbaar door $3$
  \end{center}

  zegt precies het tegenovergestelde.

  We noemen dit de \textbf{negatie} van $p$.


  \begin{definitie}
  De \textbf{negatie} van een uitspraak $p$ is de uitspraak die beweert
  dat $p$ niet waar is.

  We noteren:

  \[
  \neg p.
  \]

  We lezen dit als

  \begin{center}
  ``niet $p$''
  \end{center}

  of

  \begin{center}
  ``het is niet zo dat $p$''.
  \end{center}
  \end{definitie}


  \begin{voorbeeld}

  Neem

  \[
  p:\quad 17\text{ is een priemgetal}.
  \]

  Dan is

  \[
  \neg p:\quad 17\text{ is geen priemgetal}.
  \]

  Omdat $p$ waar is, is $\neg p$ onwaar.
  \end{voorbeeld}


  \begin{voorbeeld}

  Neem

  \[
  q:\quad 10>15.
  \]

  Dan is

  \[
  \neg q:\quad 10\leq15.
  \]

  De uitspraak $q$ is onwaar, terwijl $\neg q$ waar is.
  \end{voorbeeld}


  Dit leidt tot een eerste eenvoudige waarheidstabel.

  \[
  \begin{array}{c|c}
  p & \neg p\\
  \hline
  \text{waar} & \text{onwaar}\\
  \text{onwaar} & \text{waar}
  \end{array}
  \]


  \begin{opmerking}
  De negatie van een uitspraak heeft altijd de tegengestelde
  waarheidswaarde.

  \[
  \important{
  p\text{ waar}
  \iff
  \neg p\text{ onwaar}
  }
  \]

  en

  \[
  \important{
  p\text{ onwaar}
  \iff
  \neg p\text{ waar}.
  }
  \]
  \end{opmerking}


  % ------------------------------------------------------------
  \FVDsubsection{Correct negeren}
  % ------------------------------------------------------------

  Een negatie vormen betekent niet noodzakelijk dat we gewoon het woord
  ``niet'' ergens in de zin toevoegen.

  We moeten de betekenis van de volledige uitspraak ontkennen.


  \begin{voorbeeld}

  Beschouw

  \[
  p:\quad x>5.
  \]

  De negatie is niet

  \[
  x<5,
  \]

  want dan zouden we $x=5$ vergeten.

  De correcte negatie is

  \[
  \neg p:\quad x\leq5.
  \]
  \end{voorbeeld}


  \begin{voorbeeld}

  Beschouw

  \[
  p:\quad x=4.
  \]

  Dan is

  \[
  \neg p:\quad x\neq4.
  \]
  \end{voorbeeld}


  \begin{voorbeeld}

  Beschouw

  \[
  p:\quad n\text{ is even}.
  \]

  Als $n\in\mathbb{Z}$, dan kunnen we schrijven

  \[
  \neg p:\quad n\text{ is oneven}.
  \]

  Hier gebruiken we de eigenschap dat elk geheel getal ofwel even,
  ofwel oneven is.
  \end{voorbeeld}


  \begin{oefening}[Negaties herkennen]{\oefB \ESS}

  Geef telkens de correcte negatie.

  \begin{question}
  \[
  x<7
  \]

  \begin{oplossing}
  \[
  x\geq7.
  \]
  \end{oplossing}
  \end{question}


  \begin{question}
  \[
  x\geq2
  \]

  \begin{oplossing}
  \[
  x<2.
  \]
  \end{oplossing}
  \end{question}


  \begin{question}
  \[
  n\text{ is deelbaar door }5.
  \]

  \begin{oplossing}
  \[
  n\text{ is niet deelbaar door }5.
  \]
  \end{oplossing}
  \end{question}

  \end{oefening}


  % ============================================================
  \FVDsection{De conjunctie: en}
  % ============================================================

  Beschouw opnieuw twee uitspraken:

  \[
  p:\quad 12\text{ is even}
  \]

  en

  \[
  q:\quad 12\text{ is deelbaar door }3.
  \]

  We kunnen ze samenvoegen tot

  \begin{center}
  $12$ is even \textbf{en} $12$ is deelbaar door $3$.
  \end{center}

  Symbolisch schrijven we

  \[
  p\land q.
  \]


  \begin{definitie}
  De \textbf{conjunctie} van twee uitspraken $p$ en $q$ is de samengestelde
  uitspraak

  \[
  p\land q.
  \]

  We lezen dit als

  \begin{center}
  ``$p$ en $q$''.
  \end{center}

  De conjunctie is waar als \textbf{beide} uitspraken waar zijn.
  \end{definitie}


  Laten we alle mogelijkheden onderzoeken.

  \[
  \begin{array}{c|c|c}
  p & q & p\land q\\
  \hline
  \text{waar} & \text{waar} & \text{waar}\\
  \text{waar} & \text{onwaar} & \text{onwaar}\\
  \text{onwaar} & \text{waar} & \text{onwaar}\\
  \text{onwaar} & \text{onwaar} & \text{onwaar}
  \end{array}
  \]


  De waarheidstabel toont dat slechts één situatie de conjunctie waar maakt:

  \[
  \important{
  p\land q\text{ is waar als zowel }p\text{ als }q\text{ waar zijn.}
  }
  \]


  \begin{voorbeeld}

  Neem

  \[
  p:\quad 10\text{ is even}
  \]

  en

  \[
  q:\quad 10>7.
  \]

  Beide zijn waar.

  Dus

  \[
  p\land q
  \]

  is waar.
  \end{voorbeeld}


  \begin{voorbeeld}

  Neem

  \[
  p:\quad 9\text{ is oneven}
  \]

  en

  \[
  q:\quad 9\text{ is een priemgetal}.
  \]

  Dan is $p$ waar maar $q$ onwaar.

  Daarom is

  \[
  p\land q
  \]

  onwaar.
  \end{voorbeeld}


  % ------------------------------------------------------------
  \FVDsubsection{Een conjunctie analyseren}
  % ------------------------------------------------------------

  Een conjunctie bevat twee voorwaarden die tegelijk vervuld moeten zijn.


  \begin{voorbeeld}

  Beschouw

  \[
  x>2\land x<7.
  \]

  Deze samengestelde uitspraak betekent dat beide voorwaarden tegelijk
  moeten gelden.

  We kunnen dit ook schrijven als

  \[
  2<x<7.
  \]

  Voor $x=5$ is de uitspraak waar.

  Voor $x=8$ is de eerste voorwaarde wel waar, maar de tweede niet.
  De conjunctie is dan onwaar.
  \end{voorbeeld}


  \begin{voorbeeld}[Een domein beschrijven]

  De voorwaarde

  \[
  x\geq0\land x\leq1
  \]

  beschrijft precies de reële getallen in het interval

  \[
  [0,1].
  \]

  Logische operatoren kunnen dus gebruikt worden om verzamelingen
  of voorwaarden nauwkeurig te beschrijven.
  \end{voorbeeld}


  \begin{oefening}[Conjuncties]{\oefB \ESS}

  Gegeven:

  \[
  p:\quad 6\text{ is even},
  \]

  \[
  q:\quad 6\text{ is een priemgetal}.
  \]

  \begin{question}
  Bepaal de waarheidswaarde van $p$.

  \begin{oplossing}
  $p$ is waar.
  \end{oplossing}
  \end{question}

  \begin{question}
  Bepaal de waarheidswaarde van $q$.

  \begin{oplossing}
  $q$ is onwaar, want
  \[
  6=2\cdot3.
  \]
  \end{oplossing}
  \end{question}

  \begin{question}
  Bepaal de waarheidswaarde van

  \[
  p\land q.
  \]

  \begin{oplossing}
  De conjunctie is onwaar, omdat niet beide uitspraken waar zijn.
  \end{oplossing}
  \end{question}

  \end{oefening}


  % ============================================================
  \FVDsection{De disjunctie: of}
  % ============================================================

  Nu verbinden we twee uitspraken met het woord \textbf{of}.

  Beschouw

  \[
  p:\quad 12\text{ is even}
  \]

  en

  \[
  q:\quad 12\text{ is deelbaar door }5.
  \]

  De samengestelde uitspraak

  \begin{center}
  $12$ is even \textbf{of} $12$ is deelbaar door $5$
  \end{center}

  noteren we als

  \[
  p\lor q.
  \]


  \begin{definitie}
  De \textbf{disjunctie} van twee uitspraken $p$ en $q$ is de
  samengestelde uitspraak

  \[
  p\lor q.
  \]

  We lezen dit als

  \begin{center}
  ``$p$ of $q$''.
  \end{center}

  De disjunctie is waar zodra \textbf{minstens één} van beide uitspraken
  waar is.
  \end{definitie}


  De volledige waarheidstabel is

  \[
  \begin{array}{c|c|c}
  p & q & p\lor q\\
  \hline
  \text{waar} & \text{waar} & \text{waar}\\
  \text{waar} & \text{onwaar} & \text{waar}\\
  \text{onwaar} & \text{waar} & \text{waar}\\
  \text{onwaar} & \text{onwaar} & \text{onwaar}
  \end{array}
  \]


  Dus:

  \[
  \important{
  p\lor q\text{ is alleen onwaar als zowel }p\text{ als }q\text{ onwaar zijn.}
  }
  \]


  % ============================================================
  \FVDsection{Het wiskundige ``of'' is inclusief}
  % ============================================================

  In het dagelijkse taalgebruik heeft het woord \emph{of} niet altijd
  dezelfde betekenis.

  Iemand vraagt bijvoorbeeld:

  \begin{center}
  ``Wil je koffie of thee?''
  \end{center}

  Vaak wordt daarmee bedoeld dat je één van beide kiest.

  In de wiskundige logica betekent

  \[
  p\lor q
  \]

  echter:

  \begin{center}
  $p$ is waar, of $q$ is waar, \textbf{of beide zijn waar}.
  \end{center}


  Dit noemen we het \textbf{inclusieve of}.


  \begin{voorbeeld}

  Neem

  \[
  p:\quad 12\text{ is even}
  \]

  en

  \[
  q:\quad 12\text{ is deelbaar door }3.
  \]

  Beide uitspraken zijn waar.

  Daarom is ook

  \[
  p\lor q
  \]

  waar.
  \end{voorbeeld}


  \begin{opmerking}
  In de wiskundige logica betekent

  \[
  p\lor q
  \]

  standaard het \textbf{inclusieve} of.

  De situatie waarbij precies één van beide uitspraken waar mag zijn,
  wordt een \textbf{exclusieve disjunctie} genoemd.

  Die gebruiken we voorlopig niet als afzonderlijke logische operator.
  \end{opmerking}


  % ------------------------------------------------------------
  \FVDsubsection{Een belangrijk verschil}
  % ------------------------------------------------------------

  Vergelijk de uitspraken

  \[
  p\land q
  \]

  en

  \[
  p\lor q.
  \]

  Bij een conjunctie moeten \textbf{beide} onderdelen waar zijn.

  Bij een disjunctie volstaat het dat \textbf{minstens één} onderdeel
  waar is.

  We kunnen dat samenvatten:

  \[
  \important{
  \begin{array}{c}
  p\land q
  \quad\Longrightarrow\quad
  \text{beide voorwaarden moeten gelden}\\[2mm]
  p\lor q
  \quad\Longrightarrow\quad
  \text{minstens één voorwaarde moet gelden}
  \end{array}
  }
  \]


  \begin{oefening}[En of of?]{\oefB \ESS}

  Beschouw

  \[
  p:\quad 8>3
  \]

  en

  \[
  q:\quad 8<5.
  \]

  \begin{question}
  Bepaal de waarheidswaarde van

  \[
  p\land q.
  \]

  \begin{oplossing}
  $p$ is waar en $q$ is onwaar.

  Dus
  \[
  p\land q
  \]
  is onwaar.
  \end{oplossing}
  \end{question}


  \begin{question}
  Bepaal de waarheidswaarde van

  \[
  p\lor q.
  \]

  \begin{oplossing}
  Omdat minstens één van beide uitspraken waar is, is
  \[
  p\lor q
  \]
  waar.
  \end{oplossing}
  \end{question}

  \end{oefening}


  % ============================================================
  \FVDsection{Waarheidstabellen}
  % ============================================================

  Tot nu toe hebben we voor elke operator afzonderlijk onderzocht wanneer
  een samengestelde uitspraak waar is.

  Een \textbf{waarheidstabel} verzamelt alle mogelijke combinaties van
  waarheidswaarden systematisch.


  Bij twee uitspraken $p$ en $q$ zijn er vier mogelijke combinaties:

  \[
  \begin{array}{c|c}
  p & q\\
  \hline
  \text{waar} & \text{waar}\\
  \text{waar} & \text{onwaar}\\
  \text{onwaar} & \text{waar}\\
  \text{onwaar} & \text{onwaar}
  \end{array}
  \]


  We kunnen meerdere logische operatoren in één tabel vergelijken:

  \[
  \begin{array}{c|c|c|c}
  p & q & p\land q & p\lor q\\
  \hline
  \text{waar} & \text{waar} & \text{waar} & \text{waar}\\
  \text{waar} & \text{onwaar} & \text{onwaar} & \text{waar}\\
  \text{onwaar} & \text{waar} & \text{onwaar} & \text{waar}\\
  \text{onwaar} & \text{onwaar} & \text{onwaar} & \text{onwaar}
  \end{array}
  \]


  \begin{definitie}
  Een \textbf{waarheidstabel} is een tabel waarin voor alle mogelijke
  waarheidswaarden van de elementaire uitspraken de waarheidswaarde van
  een samengestelde uitspraak wordt bepaald.
  \end{definitie}


  Waarheidstabellen zijn meer dan een rekenschema.

  Ze stellen ons in staat om de \textbf{logische structuur} van een
  uitspraak te onderzoeken zonder afhankelijk te zijn van de concrete
  inhoud van $p$ en $q$.


  % ------------------------------------------------------------
  \FVDsubsection{Een waarheidstabel stap voor stap}
  % ------------------------------------------------------------

  Beschouw

  \[
  \neg p\land q.
  \]

  We bouwen de tabel in stappen.

  Eerst noteren we alle mogelijkheden voor $p$ en $q$.

  Vervolgens bepalen we $\neg p$.

  Ten slotte bepalen we wanneer $\neg p$ en $q$ tegelijk waar zijn.

  \[
  \begin{array}{c|c|c|c}
  p & q & \neg p & \neg p\land q\\
  \hline
  \text{waar} & \text{waar}
  & \text{onwaar} & \text{onwaar}\\

  \text{waar} & \text{onwaar}
  & \text{onwaar} & \text{onwaar}\\

  \text{onwaar} & \text{waar}
  & \text{waar} & \text{waar}\\

  \text{onwaar} & \text{onwaar}
  & \text{waar} & \text{onwaar}
  \end{array}
  \]


  \begin{opmerking}
  Bij een complexe logische uitdrukking werk je best
  \textbf{van binnen naar buiten}.

  Voeg voor tussenstappen extra kolommen toe.

  Zo verklein je de kans op fouten.
  \end{opmerking}


  \begin{oefening}[Een waarheidstabel opbouwen]{\oefB \ESS}

  Maak de waarheidstabel van

  \[
  p\land\neg q.
  \]

  \begin{oplossing}

  \[
  \begin{array}{c|c|c|c}
  p&q&\neg q&p\land\neg q\\
  \hline
  \text{waar}&\text{waar}&\text{onwaar}&\text{onwaar}\\
  \text{waar}&\text{onwaar}&\text{waar}&\text{waar}\\
  \text{onwaar}&\text{waar}&\text{onwaar}&\text{onwaar}\\
  \text{onwaar}&\text{onwaar}&\text{waar}&\text{onwaar}
  \end{array}
  \]

  De samengestelde uitspraak is dus alleen waar wanneer
  $p$ waar en $q$ onwaar is.
  \end{oplossing}

  \end{oefening}


  % ============================================================
  \FVDsection{Van gewone taal naar logische symbolen}
  % ============================================================

  Logische symbolen zijn geen doel op zich.

  Ze geven ons een compacte en precieze manier om de structuur van
  uitspraken zichtbaar te maken.


  \begin{voorbeeld}

  Gegeven:

  \[
  p:\quad n\text{ is even}
  \]

  en

  \[
  q:\quad n>10.
  \]

  De uitspraak

  \begin{center}
  $n$ is even en groter dan $10$
  \end{center}

  kan worden geschreven als

  \[
  p\land q.
  \]
  \end{voorbeeld}


  \begin{voorbeeld}

  Met dezelfde definities betekent

  \[
  p\lor q
  \]

  de uitspraak

  \begin{center}
  $n$ is even of $n$ is groter dan $10$.
  \end{center}

  Omdat het wiskundige ``of'' inclusief is, mag $n$ ook aan beide
  voorwaarden voldoen.
  \end{voorbeeld}


  \begin{voorbeeld}

  De uitspraak

  \begin{center}
  $n$ is niet even
  \end{center}

  wordt

  \[
  \neg p.
  \]
  \end{voorbeeld}


  % ------------------------------------------------------------
  \FVDsubsection{De hoofdoperator herkennen}
  % ------------------------------------------------------------

  Wanneer verschillende operatoren voorkomen, is het belangrijk om te
  zien welke bewerking de volledige uitspraak op het hoogste niveau
  opbouwt.

  Beschouw

  \[
  \neg p\land q.
  \]

  Hier vormen we eerst $\neg p$.

  Daarna verbinden we $\neg p$ en $q$ met \emph{en}.

  De hoofdoperator is dus

  \[
  \land.
  \]


  Beschouw daarentegen

  \[
  \neg(p\land q).
  \]

  Hier wordt eerst de conjunctie

  \[
  p\land q
  \]

  gevormd.

  Daarna wordt de volledige conjunctie ontkend.

  De hoofdoperator is nu

  \[
  \neg.
  \]


  \begin{opmerking}
  Haakjes zijn in logische uitdrukkingen essentieel.

  De uitdrukkingen

  \[
  \neg p\land q
  \]

  en

  \[
  \neg(p\land q)
  \]

  hebben niet dezelfde structuur en hoeven dus ook niet dezelfde
  waarheidswaarde te hebben.
  \end{opmerking}


  % ============================================================
  \FVDsection{Haakjes veranderen de betekenis}
  % ============================================================

  Beschouw bijvoorbeeld

  \[
  p=\text{waar},
  \qquad
  q=\text{onwaar}.
  \]

  Dan is

  \[
  \neg p\land q
  \]

  gelijk aan

  \[
  \text{onwaar}\land\text{onwaar},
  \]

  dus onwaar.

  Maar

  \[
  \neg(p\land q)
  \]

  wordt

  \[
  \neg(\text{waar}\land\text{onwaar}).
  \]

  Omdat

  \[
  \text{waar}\land\text{onwaar}
  \]

  onwaar is, krijgen we

  \[
  \neg(\text{onwaar}),
  \]

  dus waar.

  De twee uitdrukkingen hebben hier verschillende waarheidswaarden.


  \begin{oefening}[Structuur analyseren]{\oefU \ESS}

  Gegeven zijn

  \[
  p=\text{onwaar},
  \qquad
  q=\text{waar}.
  \]

  Bepaal de waarheidswaarde van de volgende uitspraken.

  \begin{question}
  \[
  \neg p\land q
  \]

  \begin{oplossing}
  Omdat $p$ onwaar is, is $\neg p$ waar.

  Dus
  \[
  \neg p\land q
  =
  \text{waar}\land\text{waar},
  \]
  en de uitspraak is waar.
  \end{oplossing}
  \end{question}


  \begin{question}
  \[
  \neg(p\land q)
  \]

  \begin{oplossing}
  Eerst:
  \[
  p\land q
  =
  \text{onwaar}\land\text{waar}
  =
  \text{onwaar}.
  \]

  Dus:
  \[
  \neg(p\land q)
  =
  \text{waar}.
  \]
  \end{oplossing}
  \end{question}


  \begin{question}
  Zijn de twee uitdrukkingen daarom altijd logisch gelijkwaardig?

  \begin{oplossing}
  Nee.

  Dat beide uitdrukkingen in dit ene geval dezelfde waarheidswaarde
  hebben, volstaat niet om te besluiten dat ze voor elke mogelijke
  combinatie van $p$ en $q$ dezelfde waarheidswaarde hebben.

  Om dat te onderzoeken moeten we alle gevallen vergelijken,
  bijvoorbeeld met een waarheidstabel.
  \end{oplossing}
  \end{question}

  \end{oefening}


  % ============================================================
  \FVDsection{Een eerste ontdekking met waarheidstabellen}
  % ============================================================

  Waarheidstabellen laten ons niet alleen de waarheidswaarde van één
  uitdrukking bepalen.

  We kunnen er ook twee logische uitdrukkingen mee vergelijken.


  Beschouw

  \[
  \neg(p\land q)
  \]

  en

  \[
  \neg p\lor\neg q.
  \]

  We maken een waarheidstabel.

  \[
  \begin{array}{c|c|c|c|c|c|c}
  p&q&p\land q&\neg(p\land q)&\neg p&\neg q&
  \neg p\lor\neg q\\
  \hline
  W&W&W&O&O&O&O\\
  W&O&O&W&O&W&W\\
  O&W&O&W&W&O&W\\
  O&O&O&W&W&W&W
  \end{array}
  \]

  We zien iets opvallends.

  De kolommen van

  \[
  \neg(p\land q)
  \]

  en

  \[
  \neg p\lor\neg q
  \]

  zijn identiek.

  De twee uitdrukkingen hebben dus voor elke mogelijke combinatie van
  $p$ en $q$ dezelfde waarheidswaarde.


  \begin{opmerking}
  Dit is geen toevalligheid.

  We hebben hier een belangrijke logische wet ontdekt.

  In een later hoofdstuk zullen we zulke \textbf{logische equivalenties}
  systematisch bestuderen en leren gebruiken.

  We zullen dan onder andere zien dat

  \[
  \neg(p\land q)
  \]

  dezelfde logische betekenis heeft als

  \[
  \neg p\lor\neg q.
  \]
  \end{opmerking}


  Deze ontdekking is een mooi voorbeeld van de kracht van een
  waarheidstabel:

  \[
  \important{
  \text{we onderzoeken niet één voorbeeld, maar alle mogelijke
  waarheidswaarden tegelijk.}
  }
  \]


  % ============================================================
  \FVDsection{Voorwaarden in wiskunde en wetenschappen}
  % ============================================================

  Conjuncties en disjuncties komen voortdurend voor in wiskundige en
  wetenschappelijke modellen.


  \begin{voorbeeld}[Temperatuurbewaking]

  Een machine werkt veilig wanneer de temperatuur minstens
  $20\,{}^\circ\mathrm C$ en hoogstens $70\,{}^\circ\mathrm C$ is.

  We kunnen de veiligheidsvoorwaarde schrijven als

  \[
  T\geq20\land T\leq70.
  \]

  Beide voorwaarden moeten tegelijk gelden.
  \end{voorbeeld}


  \begin{voorbeeld}[Een alarmsysteem]

  Een alarmsysteem wordt geactiveerd wanneer

  \begin{itemize}
    \item de temperatuur hoger is dan $80\,{}^\circ\mathrm C$,
    \item of de druk hoger is dan $5\,\mathrm{bar}$.
  \end{itemize}

  Definieer

  \[
  p:\quad T>80
  \]

  en

  \[
  q:\quad P>5.
  \]

  De alarmvoorwaarde wordt

  \[
  p\lor q.
  \]

  Het alarm gaat dus af wanneer één van beide voorwaarden waar is,
  maar ook wanneer beide tegelijk waar zijn.
  \end{voorbeeld}


  \begin{voorbeeld}[Selectie van data]

  Een onderzoeker wil uit een databestand alleen metingen selecteren
  waarvoor

  \[
  10\leq t\leq20.
  \]

  Logisch kunnen we dit schrijven als

  \[
  t\geq10\land t\leq20.
  \]

  Ook databanken, zoekmachines en programmeertalen gebruiken voortdurend
  zulke logische voorwaarden.
  \end{voorbeeld}


  % ============================================================
  \FVDsection{Logische voorwaarden in informatica}
  % ============================================================

  In een computerprogramma kunnen beslissingen rechtstreeks gebaseerd
  zijn op logische voorwaarden.

  Bijvoorbeeld:

  \[
  T>80
  \]

  kan waar of onwaar zijn.

  Maar ook complexere voorwaarden zijn mogelijk:

  \[
  T>80\lor P>5
  \]

  of

  \[
  T>20\land T<70.
  \]


  Een programma kan op basis van de waarheidswaarde vervolgens een
  actie uitvoeren.


  \begin{voorbeeld}[Toegangscontrole]

  Een systeem geeft toegang wanneer iemand

  \begin{itemize}
    \item een geldige badge heeft;
    \item én een correcte pincode invoert.
  \end{itemize}

  Definieer

  \[
  p:\quad \text{de badge is geldig}
  \]

  en

  \[
  q:\quad \text{de pincode is correct}.
  \]

  De toegangsvoorwaarde is

  \[
  p\land q.
  \]

  Als één van beide voorwaarden niet voldaan is, wordt de toegang
  geweigerd.
  \end{voorbeeld}


  \begin{opmerking}
  Logische operatoren vormen een fundamenteel onderdeel van onder andere

  \begin{itemize}
    \item programmeertalen;
    \item algoritmen;
    \item zoekopdrachten;
    \item digitale elektronica;
    \item regelsystemen;
    \item artificiële intelligentie;
    \item databanken.
  \end{itemize}

  Dezelfde formele logica die we in de wiskunde gebruiken, vormt dus
  ook een belangrijke basis van digitale technologie.
  \end{opmerking}


  % ============================================================
  \FVDsection{Analyseren in plaats van gokken}
  % ============================================================

  Bij eenvoudige uitspraken kunnen we vaak onmiddellijk zien of een
  samengestelde uitspraak waar is.

  Bij ingewikkeldere structuren wordt intuïtie minder betrouwbaar.


  Beschouw bijvoorbeeld

  \[
  \neg(p\lor q)\land r.
  \]

  Om deze uitdrukking correct te onderzoeken, moeten we de structuur
  herkennen.

  We kunnen ze stap voor stap ontleden:

  \[
  p\lor q
  \]

  vervolgens

  \[
  \neg(p\lor q)
  \]

  en ten slotte

  \[
  \neg(p\lor q)\land r.
  \]


  \begin{opmerking}
  Bij het analyseren van een samengestelde uitspraak zijn drie vragen
  bijzonder nuttig:

  \[
  \important{
  \begin{array}{l}
  1.\ \text{Uit welke elementaire uitspraken bestaat ze?}\\[1mm]
  2.\ \text{Welke logische operatoren komen voor?}\\[1mm]
  3.\ \text{Welke operator bepaalt de structuur van de volledige uitspraak?}
  \end{array}
  }
  \]
  \end{opmerking}


  \begin{oefening}[Ontleed de uitspraak]{\oefU \ESS}

  Beschouw

  \[
  \neg p\lor(q\land r).
  \]

  \begin{question}
  Welke elementaire uitspraken komen voor?

  \begin{oplossing}
  $p$, $q$ en $r$.
  \end{oplossing}
  \end{question}


  \begin{question}
  Welke logische operatoren komen voor?

  \begin{oplossing}
  Negatie $\neg$, disjunctie $\lor$ en conjunctie $\land$.
  \end{oplossing}
  \end{question}


  \begin{question}
  Welke bewerking gebeurt eerst binnen de haakjes?

  \begin{oplossing}
  De conjunctie
  \[
  q\land r.
  \]
  \end{oplossing}
  \end{question}


  \begin{question}
  Wat is de hoofdoperator van de volledige uitspraak?

  \begin{oplossing}
  De hoofdoperator is
  \[
  \lor.
  \]

  De volledige uitspraak verbindt $\neg p$ met $(q\land r)$ door middel
  van een disjunctie.
  \end{oplossing}
  \end{question}

  \end{oefening}



  \end{document}